% from ega0-1-fr.tex
\oldpage[0\textsubscript{I}]{24}

\emph{le couple formé de $\mathcal{F}(U)$ et de la famille $(\rho_\alpha)$ est
solution du problème universel pour $(\mathcal{F}(U_\alpha))$,
$(\mathcal{F}(U_\alpha\cap U_\beta))$ et $(\rho_{\alpha\beta})$ (3.1.1)}%
\footnote{C'est un cas particulier de la notion générale de \emph{limite
projective} (non filtrante) (\emph{voir} (T, I, 1.8) et le livre en préparation
annoncé dans l'Introduction).}.

Il revient au même de dire que, pour tout objet $T$ de $\mathbf{K}$, la
famille $U\to\operatorname{Hom}(T,\mathcal{F}(U))$ est un \emph{faisceau
d'ensembles}.
\end{env}

\begin{env}[3.1.3]
\label{0.3.1.3-fr}
Supposons que $\mathbf{K}$ soit la catégorie définie par une «~espèce de
structure avec morphismes~» $\Sigma$, les objets de $\mathbf{K}$ étant donc
les ensembles munis de structures d'espèce $\Sigma$ et les morphismes ceux de
$\Sigma$. Supposons que la catégorie $\mathbf{K}$ vérifie en outre la
condition suivante~:

\smallskip
\noindent
\textup{(E)} Si $(A,(\rho_\alpha))$ est solution d'un problème d'application
universelle \emph{dans la catégorie $\mathbf{K}$} pour des familles
$(A_\alpha)$, $(A_{\alpha\beta})$, $(\rho_{\alpha\beta})$, alors c'est aussi
une solution du problème d'application universelle pour les mêmes familles
\emph{dans la catégorie des ensembles} (c'est-à-dire quand on considère $A$,
les $A_\alpha$ et $A_{\alpha\beta}$ comme des ensembles, les $\rho_\alpha$ et
$\rho_{\alpha\beta}$ comme des applications)%
\footnote{On peut prouver que cela signifie aussi que le foncteur canonique
$\mathbf{K}\to(\mathrm{Ens})$ \emph{permute aux limites projectives} (non
nécessairement filtrantes).}.

Dans ces conditions, la condition (F) entraîne que, considéré comme
préfaisceau \emph{d'ensembles}, $U\to\mathcal{F}(U)$ est un \emph{faisceau}.
En outre, pour qu'une application $u:T\to\mathcal{F}(U)$ soit un morphisme de
$\mathbf{K}$, il faut et il suffit, en vertu de (F), que chaque application
$\rho_\alpha\circ u$ soit un morphisme $T\to\mathcal{F}(U_\alpha)$, ce qui
signifie que la structure d'espèce $\Sigma$ sur $\mathcal{F}(U)$ est
\emph{structure initiale} pour les morphismes $\rho_\alpha$. Réciproquement,
supposons qu'un préfaisceau $U\to\mathcal{F}(U)$ sur $X$, à valeurs dans
$\mathbf{K}$, soit un \emph{faisceau d'ensembles}, et vérifie la condition
précédente~; il est clair alors qu'il satisfait à (F), donc est un
\emph{faisceau à valeurs dans $\mathbf{K}$}.
\end{env}

\begin{env}[3.1.4]
\label{0.3.1.4-fr}
Lorsque $\Sigma$ est l'espèce de structure de groupe ou d'anneau, le fait que
le préfaisceau $U\to\mathcal{F}(U)$ à valeurs dans $\mathbf{K}$ est un
\emph{faisceau d'ensembles} entraîne \emph{ipso facto} que c'est un
\emph{faisceau à valeurs dans $\mathbf{K}$} (autrement dit, un faisceau de
groupes ou d'anneaux au sens de (G))%
\footnote{Cela tient à ce que dans la catégorie $\mathbf{K}$, tout morphisme
qui est une \emph{bijection} (en tant qu'application d'ensembles) est un
\emph{isomorphisme}. Cela n'est plus vrai lorsque $\mathbf{K}$ est la
catégorie des espaces topologiques, par exemple.}. Mais il n'en est plus de
même lorsque par exemple $\mathbf{K}$ est la catégorie des \emph{anneaux
topologiques} (avec pour morphismes les représentations continues)~: un
faisceau à valeurs dans $\mathbf{K}$ est un faisceau d'anneaux
$U\to\mathcal{F}(U)$ tel que, pour tout ouvert $U$ et tout recouvrement de $U$
par des ouverts $U_\alpha\subset U$, la topologie de l'anneau $\mathcal{F}(U)$
soit \emph{la moins fine} rendant continues les représentations
$\mathcal{F}(U)\to\mathcal{F}(U_\alpha)$. On dira dans ce cas que
$U\to\mathcal{F}(U)$ considéré comme faisceau d'anneaux (sans topologie) est
\emph{sous-jacent} au faisceau d'anneaux topologiques
$U\to\mathcal{F}(U)$. Les morphismes
$u_V:\mathcal{F}(V)\to\mathcal{G}(V)$ ($V$ ouvert arbitraire de $X$) de
faisceaux d'anneaux topologiques sont donc des homomorphismes des faisceaux
d'anneaux sous-jacents, tels que $u_V$ soit \emph{continu} pour tout ouvert
$V\subset X$~; pour les distinguer des homomorphismes quelconques des
faisceaux d'anneaux sous-jacents, on les appellera homomorphismes
\emph{continus} de faisceaux d'anneaux topologiques. On a des définitions et
conventions analogues pour les faisceaux d'espaces topologiques ou de groupes
topologiques.
\end{env}

